% Part: first-order-logic
% Chapter: tableaux
% Section: identity

\documentclass[../../../include/open-logic-section]{subfiles}

\begin{document}

\olfileid{fol}{tab}{ide}

\olsection{له \usetoken{S}{identity} سره \usetoken{P}{tableau}}

له !!{identity} سره !!^{tableau}s د استنتاج اضافي قاعدې غواړي۔
د~$\eq$ قاعدې دا دي ($t$، $t_1$ او $t_2$ تړلي ترمونه دي):

\begin{defish}
\AxiomC{}
\RightLabel{$\eq$}
\UnaryInfC{\sFmla{\True}{\eq[t][t]}}
\DisplayProof
\hfill
\AxiomC{\sFmla{\True}{\eq[t_1][t_2]}}
\noLine
\UnaryInfC{\sFmla{\True}{!A(t_1)}}
\RightLabel{$\TRule{\True}{\eq}$}
\UnaryInfC{\sFmla{\True}{!A(t_2)}}
\DisplayProof
\hfill
\AxiomC{\sFmla{\True}{\eq[t_1][t_2]}}
\noLine
\UnaryInfC{\sFmla{\False}{!A(t_1)}}
\RightLabel{$\TRule{\False}{\eq}$}
\UnaryInfC{\sFmla{\False}{!A(t_2)}}
\DisplayProof
\end{defish}
پام وکړئ چې له نورو ټولو قاعدو سره په توپير کښې،
$\TRule{\True}{\eq}$ او $\TRule{\False}{\eq}$ غواړي چې \emph{دوه}
نښه لرونکي !!{formula}s له وړاندې پر څانګه وي؛ يعنې دواړه
$\sFmla{\True}{\eq[t_1][t_2]}$ او $\sFmla{S}{!A(t_1)}$۔

\begin{ex}
که $s$ او $t$ تړلي ترمونه وي، نو
$\eq[s][t], !A(s) \Proves !A(t)$:
\begin{oltableau}
  [\sFmla{\False}{\formula{A}(t)}, just = \TAss
    [\sFmla{\True}{\eq[s][t]}, just = \TAss
      [\sFmla{\True}{\formula{A}(s)}, just = \TAss
        [\sFmla{\True}{\formula{A}(t)}, just={\TRule{\True}{\eq}[2, 3]}, close]
      ]
    ]
  ]
\end{oltableau}
کېدے شي دا د يو شان څيزونو د تعويض د اصل يا د لايبنېڅ د قانون په
نوم درته اشنا وي۔

!!^{tableau}s ثابتوي چې $\eq$ متناظره ده، يعنې
$\eq[s_1][s_2] \Proves \eq[s_2][s_1]$:
\begin{oltableau}
  [\sFmla{\False}{\eq[s_2][s_1]}, just = \TAss
    [\sFmla{\True}{\eq[s_1][s_2]}, just = \TAss
      [\sFmla{\True}{\eq[s_1][s_1]}, just = {$\eq$}
        [\sFmla{\True}{\eq[s_2][s_1]}, just = {\TRule{\True}{\eq}[2, 3]}, close]
      ]
    ]
  ]
\end{oltableau}
دلته کرښه~$2$ د~$\TRule{\True}{\eq}$ لومړنے اړين
!!{formula}~$\sFmla{\True}{\eq[s_1][s_2]}$ دے۔ کرښه~$3$ دويم
اړين فارمول دے، چې د~$\sFmla{\True}{!A(s_1)}$ بڼه لري---$!A(x)$
د~$\eq[x][s_1]$ په توګه وګڼئ؛ بيا $!A(s_1)$،
$\eq[s_1][s_1]$ او $!A(s_2)$، $\eq[s_2][s_1]$ کېږي۔
\textbf{د ژباړې د نسخې سمون (OLTAB-010):} سرچينې کرښه درې د اې
د دويم ترم بڼه بللې وه، خو د قاعدې په دې تطبيق کښې هغه د اې د
لومړي ترم بڼه ده، لکه ورپسې تشريح چې هم ښيي۔

تابلوګان دا هم ثابتوي چې $\eq$ انتقالي ده، يعنې
$\eq[s_1][s_2], \eq[s_2][s_3] \Proves \eq[s_1][s_3]$:
\begin{oltableau}
  [\sFmla{\False}{\eq[s_1][s_3]}, just = \TAss
    [\sFmla{\True}{\eq[s_1][s_2]}, just = \TAss
      [\sFmla{\True}{\eq[s_2][s_3]}, just = \TAss
        [\sFmla{\True}{\eq[s_1][s_3]}, just = {\TRule{\True}{\eq}[3, 2]}, close]
      ]
    ]
  ]
\end{oltableau}
په دې !!{tableau} کښې د~$\TRule{\True}{\eq}$ لومړنے اړين
!!{formula} کرښه~$3$، $\sFmla{\True}{\eq[s_2][s_3]}$ ده
($s_2$ د~$t_1$ او $s_3$ د~$t_2$ ونډه لري)۔ دويم اړين فارمول چې
د~$\sFmla{\True}{!A(s_2)}$ بڼه لري، کرښه~$2$ ده۔ دلته $!A(x)$
د~$\eq[s_1][x]$ په توګه وګڼئ؛ په دې سره $!A(s_2)$،
$\eq[s_1][s_2]$ (يعنې کرښه~$2$) او $!A(s_3)$ په پايله کښې
!!{formula}~$\eq[s_1][s_3]$ کېږي۔
\textbf{د ژباړې د نسخې سمون (OLTAB-011):} سرچينې د اې د دويم
ترم د بڼې د پراخولو پر ځاے د قاعدې عمومي برابري ليکلې وه؛ خو
د همدلته ټاکل شوي اې له مخې اړوند فارمول د لومړي او دويم ترم
برابري ده۔
\end{ex}

\begin{prob}
د لاندې دپاره تړلي !!{tableau}s ورکړئ:
\begin{enumerate}
\item $\sFmla{\False}{\lforall[x][\lforall[y][((x = y \land !A(x))
      \lif !A(y))]]}$
\item $\sFmla{\False}{\lexists[x][(!A(x) \land
    \lforall[y][(!A(y) \lif y = x)])]}$,\\
  $\sFmla{\True}{\lexists[x][!A(x)] \land
    \lforall[y][\lforall[z][((!A(y) \land !A(z)) \lif y = z)]]}$
\end{enumerate}
\end{prob}

\end{document}
